\begin{beginnerstatement}[In simple terms: Generating a project]

The \texttt{init} command starts the generator for a new product project. The
generator builds a \emph{scaffold} from the selected profile: a prepared basic
structure of folders and files. The target directory is the location where
this project will be created. A \emph{dry run} checks the intended process
without writing files. Project identity includes details such as the name and
short name that distinguish the generated project from the template. An
application ID uniquely identifies a Tauri application. A
\emph{reverse-domain identifier} is such an ID in which words, as in
\texttt{com.example.product}, are arranged from the general Internet domain
to the specific product.

\end{beginnerstatement}
